% =====================================================================
% Rozdział 7: Ewolucja, rozszerzanie i wersjonowanie modelu LEM
% =====================================================================

\chapter{Ewolucja, rozszerzanie i wersjonowanie modelu LEM}
\label{chap:ewolucja}

\section{Klasyfikacja zmian modelu}
Zmiany dotyczące modelu LEM MUSZĄ być klasyfikowane według ich wpływu na semantykę modelu. Poniższa klasyfikacja dotyczy zmian w modelu podstawowym; zasady ewolucji samych profili określa sekcja 4.6, która odwołuje się odpowiednio do niniejszej klasyfikacji.

Wyróżnia się następujące typy zmian:

\subsection{Zmiana korekcyjna}
Zmiana korekcyjna usuwa niejednoznaczności, błędy dokumentacyjne lub poprawia opis istniejących elementów modelu.

Zmiana korekcyjna NIE MOŻE zmieniać interpretacji komponentów modelu LEM ani wpływać na zgodne implementacje istniejących wersji.

\subsection{Zmiana rozszerzająca}
Zmiana rozszerzająca wprowadza nowe możliwości modelu przy zachowaniu istniejącej semantyki.

Zmiana rozszerzająca MOŻE obejmować: dodanie nowych komponentów, dodanie nowych opcjonalnych właściwości, rozszerzenie zakresu obsługiwanych reprezentacji, dodanie nowych profili zastosowań.

Zmiana rozszerzająca NIE MOŻE powodować zmiany znaczenia istniejących komponentów modelu LEM.

\subsection{Zmiana semantyczna}
Zmiana semantyczna wpływa na podstawowe znaczenie elementów modelu LEM.

Przykłady zmian semantycznych obejmują: zmianę definicji istniejącego komponentu, zmianę relacji pomiędzy komponentami, zmianę zasad interpretacji kanonicznej, zmianę zasad określania równoważności semantycznej.

Zmiana semantyczna MUSI być traktowana jako zmiana podstawowej wersji modelu LEM.

\section{Zasady wersjonowania}
Wersjonowanie LEM MUSI umożliwiać jednoznaczne określenie zgodności pomiędzy implementacjami wykorzystującymi model.

Wersjonowanie powinno rozróżniać: wersję modelu semantycznego LEM, wersję zasad reprezentacji oraz wersję profili zastosowań.

Zmiana wersji reprezentacji lub profilu NIE MOŻE powodować zmiany podstawowej semantyki modelu LEM.

\section{Kompatybilność wersji}
Implementacje LEM MUSZĄ określać wersję modelu semantycznego, według której interpretowane są dane lokalizacyjne.

Dwie implementacje są zgodne semantycznie, jeżeli:
\begin{itemize}
    \item wykorzystują zgodną wersję modelu semantycznego LEM,
    \item interpretują komponenty zgodnie z tym samym modelem znaczeniowym,
    \item zachowują zasady interpretacji kanonicznej.
\end{itemize}

Różnice w wersjach reprezentacji lub profili mogą być obsługiwane niezależnie od wersji rdzenia semantycznego, o ile zachowana jest zgodność interpretacji.

\section{Zasady rozszerzania modelu}
Rozszerzenia modelu LEM MUSZĄ zachowywać istniejącą semantykę komponentów podstawowych.

Nowy komponent lub właściwość:
\begin{itemize}
    \item NIE MOŻE zmieniać znaczenia istniejących komponentów,
    \item NIE MOŻE powodować alternatywnej interpretacji istniejących danych,
    \item MUSI posiadać jednoznacznie określone znaczenie.
\end{itemize}

Jeżeli wymagane jest wprowadzenie nowego znaczenia informacji lokalizacyjnej, zmiana MUSI zostać wykonana w modelu podstawowym LEM, a nie poprzez profil lub implementację lokalną.

Niniejsza sekcja dotyczy rozszerzeń przeznaczonych do formalnego włączenia do modelu LEM. Rozszerzenia o charakterze lokalnym, niepodlegające takiemu włączeniu, opisuje sekcja 7.5.

\section{Rozszerzenia prywatne i domenowe}
Implementacje MOGĄ definiować rozszerzenia lokalne lub domenowe, przeznaczone wyłącznie do użytku w określonym środowisku i niepodlegające formalnemu włączeniu do modelu podstawowego, jeżeli rozszerzenia te:
\begin{itemize}
    \item nie zmieniają znaczenia komponentów zdefiniowanych przez LEM,
    \item nie powodują niejednoznacznej interpretacji danych zgodnych z LEM,
    \item są wyraźnie oddzielone od podstawowego modelu LEM.
\end{itemize}

Rozszerzenia lokalne NIE MOGĄ być traktowane jako standardowe komponenty LEM bez formalnego włączenia ich do modelu zgodnie z zasadami sekcji 7.4.

\section{Proces zmian semantycznych}
Zmiany wpływające na podstawową semantykę modelu LEM MUSZĄ podlegać szczególnej procedurze wersjonowania.

Procedura ta MUSI zapewniać: identyfikację zakresu zmiany, określenie wpływu na istniejące implementacje, określenie zasad migracji oraz zachowanie możliwości interpretacji starszych danych.

Zmiana semantyczna NIE MOŻE być wprowadzona wyłącznie poprzez zmianę reprezentacji lub profilu zastosowania.

Sposób prowadzenia tego procesu w praktyce (np. organ decyzyjny, tryb konsultacji, publikacja zmian) pozostaje poza zakresem modelu LEM.

\section{Obsługa nieznanych elementów}
Implementacje LEM MUSZĄ rozróżniać elementy nieznane według ich znaczenia dla interpretacji danych. To, czy dany element ma dla implementacji charakter opcjonalny czy krytyczny, jest ustalane kontekstowo — na podstawie wymagań profilu zastosowania (zob. 4.2, 4.4.2), a nie odgórnie przez sam model LEM.

\subsection{Element opcjonalny}
Element opcjonalny, którego brak nie wpływa na interpretację podstawowej informacji lokalizacyjnej, MOŻE zostać pominięty przez implementację.

\subsection{Element krytyczny}
Element krytyczny, którego znaczenie jest wymagane do poprawnej interpretacji informacji lokalizacyjnej, NIE MOŻE zostać pominięty.

Implementacja, która nie obsługuje elementu krytycznego, MUSI oznaczyć dane jako nieinterpretowalne lub niezgodne z obsługiwaną wersją modelu (zob. 5.4).

\section{Stabilność modelu LEM}
Rozwój modelu LEM MUSI zachowywać stabilność podstawowej semantyki.

Zmiany modelu powinny preferować: rozszerzanie zamiast redefinicji, dodawanie nowych możliwości zamiast zmiany istniejących znaczeń, zachowanie kompatybilności istniejących danych.

Podstawowym celem procesu ewolucji LEM jest zachowanie możliwości jednoznacznej interpretacji informacji lokalizacyjnej w kolejnych wersjach modelu.